% ID: SEL-D0050-OA-E04
% ESTADO: cotejada. Segunda pasada visual y borrador compilado revisados.
% Fuente: Selectividad/sel_2013_matematicas_aplicadas/junio/titular/EXAMEN 6 -ANDALUCÍA.pdf
% SHA-256: 1f0983e351b8435c5b6c7b68b05fe2a64034faf77a85a1f9480407aa0f2a88aa
% Página: 1; opción A; ejercicio original 4.

Se quiere estimar la proporción de hembras entre los peces de una piscifactoría; para ello se ha tomado una muestra aleatoria de 500 peces, y en ella hay 175 hembras.
\begin{enumerate}[label=\alph*)]
\item \textbf{(1.5 puntos)} Calcule un intervalo de confianza para la proporción de hembras en esta población de peces, con un nivel de confianza del 94\%.
\item \textbf{(1 punto)} A la vista del resultado del muestreo se quiere repetir la experiencia para conseguir un intervalo de confianza con el mismo nivel y un error máximo de 0.02, \textquestiondown{}cuál es el tamaño mínimo que debe tener la nueva muestra?
\end{enumerate}
